\subsection{Image Reconstruction (CIFAR-10)}
\label{sec:reconstruction}

\paragraph{Task Description.}
The reconstruction experiment evaluates how gracefully the chord substrate
degrades as the occlusion level increases.
Each CIFAR-10 image ($32 \times 32$, decoded to raw NRGBA pixel bytes,
4\,096 bytes total) is ingested once into the unified substrate.
The same image is then queried at nine holdout levels (10\%--90\% RIGHT),
where the last $k$\% of pixel bytes are withheld and must be recovered
purely from chord attractor resonance, with no explicit image model,
convolution, or colour-space interpolation.

\paragraph{Results.}
Figure~\ref{fig:reconstruction_scaling} shows the \emph{occlusion scaling
curve}: mean scores plotted against holdout percentage.
A graceful decay (slow slope) would indicate that the chord substrate
encodes global spatial structure, whereas a cliff at high holdout
percentages is expected once the missing region exceeds the attractor's
effective support radius.

Figure~\ref{fig:reconstruction_strip} shows the qualitative result at the
representative 50\% holdout level: each row presents the Original, the
Prompt (bottom half zeroed), the Reconstructed image, and an automatically
computed Error Heatmap (green = zero error, red = maximum error).

Across all holdout levels the overall weighted score was 0.061.

\paragraph{Assessment.}
Reconstruction quality was low across all holdout levels.  At this ingestion
scale the chord substrate has not accumulated sufficient pixel-level
redundancy to anchor reliable attractor retrieval across the diverse
colour distributions of CIFAR-10 imagery.
